% Part: proof-theory
% Chapter: natural-deduction
% Section: grafting

\documentclass[../../../include/open-logic-section]{subfiles}

\begin{document}

\olfileid{pt}{nat}{gra}

\olsection{في تطعيم \usetoken{P}{proof}}

!!^{proof}s في~\Log{N1c} و\Log{N1i} إنما هي~!!{proof}s لـ!!{formula}s~$!A$
من فروض~$!B$. فإذا كان لـ$!B$ نفسها~!!a{proof}، سألنا: كيف نصل~!!{proof}
$!B$ بـ!!{proof}~$!A$ من~$!B$، حتى نحصل على~!!a{proof} لـ$!A$ لا
يستند إلى الفرض~$!B$؟

ومن الوجوه أن نحذف الفرض~$!B$ من~!!{proof}~$!A$ بإعمال
$\Intro{\lif}$، ثم نصل~!!{proof} الناتج لـ$!B \lif !A$
بـ!!{proof}~$!B$، فينتج
\[
\AxiomC{$\Discharge{!B}{{\OLId{latin-ordinary}{x}}}$}
\DeduceC{$!A$}
\DischargeRule{\Intro{\lif}}{{\OLId{latin-ordinary}{x}}}
\UnaryInfC{$!B \lif !A$}
\AxiomC{}
\DeduceC{$!B$}
\RightLabel{\Elim{\lif}}
\BinaryInfC{$!A$}
\DisplayProof
\]
وهذا طريق مباشر، غير أنه مشوب بفضول: فما وجه إدخال~$\lif$ إذا كنا
نحذفه من فورنا؟ والأولى أن نستغني عن هذا «الالتفاف» المار بـ~$!B \lif !A$.
(ومؤدى مبرهنة التطبيع أن الاستغناء عن هذه الالتفافات ممكن دائمًا.)

وفي~\Log{N1c} و\Log{N1i} يسهل اجتنابه: نقيم مقام الفروض~$!B^{\OLId{latin-ordinary}{x}}$
كامل~!!{proof}~$!B$. ولما كان «تطعيم»~!!{proof}s بفروض براهين أخرى كثير
النفع، أثبتناه في التعريف الآتي.

\begin{defn}
ليكن~$\delta_{\OLMathNumeral{1}}$~!!{proof} في~\Log{N1c} (أو~!!{proof} في
\Log{N1i}) لـ$!A$ من الفرض المفتوح~$!B^{\OLId{latin-ordinary}{x}}$، وليكن~$\delta_{\OLMathNumeral{2}}$~!!{proof}
في~\Log{N1c} (أو~!!{proof} في~\Log{N1i}) لـ$!B$. فنعرف
$\Subst{\delta_{\OLMathNumeral{1}}}{\delta_{\OLMathNumeral{2}}}{!B^{\OLId{latin-ordinary}{x}}}$ بأنه ما يحصل من إحلال~$\delta_{\OLMathNumeral{2}}$ محل
كل ظهور غير مبرأ للفرض~$!B^{\OLId{latin-ordinary}{x}}$ في~$\delta_{\OLMathNumeral{1}}$.
\end{defn}

إذا لم يكن في~$\delta_{\OLMathNumeral{1}}$ و$\delta_{\OLMathNumeral{2}}$~!!{formula}s فرض متغايرة تحمل
علامة واحدة، ولم يشتمل فرض مفتوح في~$\delta_{\OLMathNumeral{2}}$ على متغير مميز من
$\delta_{\OLMathNumeral{1}}$، كان~$\Subst{\delta_{\OLMathNumeral{1}}}{\delta_{\OLMathNumeral{2}}}{!B^{\OLId{latin-ordinary}{x}}}$~!!{proof} صحيحًا،
وفروضه المفتوحة مجموع فروض~$\delta_{\OLMathNumeral{1}}$ وفروض~$\delta_{\OLMathNumeral{2}}$. وهذان الشرطان
متحققان عند تطعيم~!!{proof}s عادة؛ فالأول حاصل متى كان~$\delta_{\OLMathNumeral{1}}$ و$\delta_{\OLMathNumeral{2}}$
برهانين جزئيين من~!!{proof} أكبر، والثاني إن تخلف أصلحناه بإعادة تسمية
المتغيرات المميزة في~$\delta_{\OLMathNumeral{1}}$.

ولهذا التعريف نظير في~\Log{N2c} و\Log{N2i}.

\begin{defn}
إذا كان~$\delta_{\OLMathNumeral{1}}$~!!{proof} في~\Log{N2c} (أو~!!{proof} في
\Log{N2i}) لـ${\OLId{latin-ordinary}{x}}:!B, {\OLId{greek-variable}{Gamma}}_{\OLMathNumeral{1}} \Sequent !A$، وكان~$\delta_{\OLMathNumeral{2}}$~!!{proof}
في~\Log{N2c} (أو~!!{proof} في~\Log{N2i}) لـ
${\OLId{greek-variable}{Gamma}}_{\OLMathNumeral{2}} \Sequent !B$، فإن~$\Subst{\delta_{\OLMathNumeral{1}}}{\delta_{\OLMathNumeral{2}}}{{\OLId{latin-ordinary}{x}}:!B}$ هو
ناتج استبدال~$\delta_{\OLMathNumeral{2}}$ بكل مسلّمة~${\OLId{latin-ordinary}{x}}:!B \Sequent !B$ في~$\delta_{\OLMathNumeral{1}}$،
وإضافة~${\OLId{greek-variable}{Gamma}}_{\OLMathNumeral{2}}$ إلى سياق كل متتالية تقع تحت تلك المسلّمات.
\end{defn}

\begin{prop}
إذا كان~$\delta_{\OLMathNumeral{1}} \Proves {\OLId{latin-ordinary}{x}}:!B, {\OLId{greek-variable}{Gamma}}_{\OLMathNumeral{1}} \Sequent !A$ و
$\delta_{\OLMathNumeral{2}} \Proves {\OLId{greek-variable}{Gamma}}_{\OLMathNumeral{2}} \Sequent !B$، فإن
$\Subst{\delta_{\OLMathNumeral{1}}}{\delta_{\OLMathNumeral{2}}}{{\OLId{latin-ordinary}{x}}:!B} \Proves {\OLId{greek-variable}{Gamma}}_{\OLMathNumeral{1}}, {\OLId{greek-variable}{Gamma}}_{\OLMathNumeral{2}} \Sequent
!A$، بشرط ألا يظهر أي متغير مميز من~$\delta_{\OLMathNumeral{1}}$ في~${\OLId{greek-variable}{Gamma}}_{\OLMathNumeral{2}}$.
\end{prop}


\begin{proof}
  بالاستقراء على~$\pheight{\delta_{\OLMathNumeral{1}}}$.
\end{proof}

\end{document}
